% =====================================================================================
% Chapter 10 · Staggered rollout — difference-in-differences
%
% Built from notebooks/08_rollout_did.ipynb. NOT ONE NUMBER IN THIS FILE IS TYPED BY HAND:
% every figure in the prose is a macro from build/macros.tex, generated by cmp.report from
% the executed notebook. The only literals below are (i) the constants of the data-generating
% model in §10.2, which are *definitions* rather than results, and (ii) the decision-rule
% thresholds (0.9 / 0.5) and the nominal 90 % interval level, which are *conventions* stated
% in the text, not measurements.
% =====================================================================================
\chapter{Staggered rollout: difference-in-differences}
\label{ch:did}

% This chapter's argument is completed by Chapter 9 (geo lift), whose Bayesian interval failed in
% the OPPOSITE direction. Chapter 9's labels are in scope in the bound book but not in this
% chapter's standalone offprint, so the cross-reference degrades to prose there rather than
% setting a "??".
\makeatletter
\newcommand{\refgeo}{\ifcsname r@ch:geo\endcsname\cref{ch:geo}\else the geo-lift chapter\fi}
\newcommand{\Refgeo}{\ifcsname r@ch:geo\endcsname\Cref{ch:geo}\else The geo-lift chapter\fi}
\makeatother

\begin{quote}\small
\emph{A loyalty app went live in \nbEightNTreated{} of a chain's \nbEightNStores{} stores in week
\nbEightLaunch, and the pilot stores' revenue rose --- but so did everyone's, because a seasonal
tide lifts all boats. Difference-in-differences nets the tide out and prices the app at
\eur{\nbEightAttCl} per store per week against a planted \eur{\nbEightTrueEffect}, which at a
running cost of \eur{\nbEightAppCost} is worth \eur{\nbEightAnnualMeanM}\,m a year across the
\nbEightNChain-store chain. Two things then happen that the folklore does not predict. Clustering
the standard error on store \emph{shrinks} it, from \eur{\nbEightAttIidSe} to
\eur{\nbEightAttClSe} --- while inflating the standard error of a different coefficient of the
same regression from \eur{\nbEightLvlIidSe} to \eur{\nbEightLvlClsSe}: the iid formula is not
conservative and not anti-conservative, it is wrong in a direction that flips inside one fit. And
the pooled Bayesian model returns an interval \nbEightBayesOverCluster$\times$ wider than either
classical arm --- an error in the \emph{opposite} direction to the geo-lift
chapter's, which is the proof that what fails is the specification and not the paradigm.}
\end{quote}

% =====================================================================================
\section{The decision}
\label{sec:did:decision}

A retail chain of \nbEightNChain{} stores has built a loyalty app and piloted it. The app went
live in \nbEightNTreated{} stores in week \nbEightLaunch; \nbEightNStores{} stores were watched in
total, for \nbEightNWeeks{} weeks, so the panel holds \nbEightNObs{} store-weeks. Revenue in the
pilot stores is up. The vendor wants a chain-wide contract at \eur{\nbEightAppCost} per store per
week, and somebody has to say yes or no.

The first thing to notice is that revenue is up \emph{everywhere}. Retail revenue moves with a
seasonal tide that hits every store in the chain, and the pilot stores were watched through the
same weeks as the others. A before-and-after comparison inside the pilot stores would therefore
book the tide as app effect and recommend the contract on evidence that contains none. That
shortcut is not a straw man: it is what a spreadsheet produces by default, and
\cref{fig:did:panel} shows exactly the wave it would be mistaking for a treatment.

The cost of that mistake is not abstract. At \eur{\nbEightAppCost} per store per week across
\nbEightNChain{} stores and \nbEightWeeksYear{} weeks, the contract the naive read would recommend
commits \eur{\nbEightContractAnnualM}\,m a year --- against an evidence base that is, on that read,
a seasonal wave. The same arithmetic runs in the other direction: if the app does work and the
chain declines it out of caution, the forgone lift is of the same order (\cref{sec:did:euro} puts
it at \eur{\nbEightAnnualMeanM}\,m a year net of the cost). Both errors are eight-figure errors, and
the only thing standing between them is whether the tide is netted out correctly.

The second thing to notice is that the pilot stores are not the control stores. Chains do not pick
pilot sites by coin flip; they pick the flagship, the tech-friendly manager, the store whose
district director volunteered. So the pilot stores may sell more --- or less --- than the controls
for a hundred reasons that have nothing to do with an app. A raw treated-minus-control comparison
of \emph{levels} books all of it as effect.

Difference-in-differences (DiD) is the workhorse that survives both objections, and it survives
them by taking two differences and subtracting them:
%
\begin{equation}
  \widehat{\text{DiD}}
  \;=\;
  \underbrace{\big(\bar y^{\,T}_{\text{post}} - \bar y^{\,T}_{\text{pre}}\big)}_{\text{app effect} \;+\; \text{tide}}
  \;-\;
  \underbrace{\big(\bar y^{\,C}_{\text{post}} - \bar y^{\,C}_{\text{pre}}\big)}_{\text{tide alone}},
  \label{eq:did:twobytwo}
\end{equation}
%
with $T$ the pilot stores and $C$ the controls. The first difference contains the app effect plus
the tide; the second is the tide by itself, measured on stores the app never touched. Subtracting
removes the tide. And because each store enters \emph{through its own change}, any time-invariant
store characteristic --- a flagship's larger footprint, a rich district, a good manager --- appears
in the store's ``before'' and in its ``after'' and cancels exactly. \textbf{The confounder does not
have to be observed, or even named. It has to be constant.} That is the whole appeal of the method,
and it is worth stating in the vocabulary of the earlier chapters: DiD does not close the back door
through the store baseline by conditioning on it. It differences it away.

What the bargain costs is stated in \cref{sec:did:ident} and is untestable. The rest of the chapter
is about paying that price with open eyes --- and about the two places where the machinery, rather
than the identification, is where the money is lost.

\input{build/floats/nb08_panel}

% =====================================================================================
\section{The data-generating model}
\label{sec:did:dgm}

A causal method cannot be graded on real data, because on real data the counterfactual is missing
--- that is the entire problem. So the estimator is first graded against a world whose answer is
known by construction. \Refgeo{} performed the complementary check on a real randomized geo
experiment; here the simulator carries the whole chapter, and it is built so that the assumption
under test holds \emph{exactly}, which is what makes the diagnostics of \cref{sec:did:valid}
interpretable as diagnostics rather than as symptoms.

Stores are indexed $s = 1, \dots, \nbEightNStores$, of which a random half ($D_s = 1$,
\nbEightNTreated{} stores) receive the app; weeks are $t = 0, \dots, \nbEightNWeeks - 1$, and the
launch is week $t_0 = \nbEightLaunch$. Observed revenue is
%
\begin{equation}
  Y_{st}
  \;=\;
  \underbrace{\alpha_s}_{\text{store baseline}}
  \;+\;
  \underbrace{60 \sin\!\Big(\tfrac{2\pi t}{12}\Big)}_{\text{shared seasonal tide}}
  \;+\;
  \underbrace{400 \, D_s \, \mathbf 1[\,t \ge 12\,]}_{\text{the app}}
  \;+\;
  \varepsilon_{st},
  \label{eq:did:dgp}
\end{equation}
%
with
%
\begin{equation}
  \alpha_s \sim \mathcal N\!\left(1000,\ 150^{2}\right),
  \qquad
  \varepsilon_{st} \sim \mathcal N\!\left(0,\ 80^{2}\right).
  \label{eq:did:dgpnoise}
\end{equation}
%
Three features of \cref{eq:did:dgp} are load-bearing, and each is planted deliberately.

\emph{The seasonal term is identical across stores.} There is no store-specific loading on it: the
tide lifts every store by the same euros in the same week. Together with a store baseline
$\alpha_s$ that does not move with $t$, this makes \textbf{parallel trends hold exactly, by
construction} --- absent the app, treated and control means move in lockstep. A chapter whose
central diagnostic is a test \emph{for} parallel trends must first be run in a world where they
genuinely hold, so that a flat pre-trend can be recognised as a true negative rather than hoped to
be one. \Cref{sec:did:wrong} then breaks each of the chapter's assumptions on purpose --- parallel
trends first, by giving the pilot stores a trend of their own; then no-anticipation; then the
estimator itself, under staggered adoption --- and watches the same diagnostic fire, or fail to.

\emph{The store baseline is large.} Its standard deviation, \eur{\nbEightAlphaSd}, is nearly twice
the weekly idiosyncratic noise, \eur{\nbEightNoiseSd}. This is realistic --- stores differ far more
from each other than a store differs from itself week to week --- and it is the single number that
drives everything surprising in this chapter. It is what the DiD contrast \emph{cancels}, what the
iid standard error \emph{charges the estimate for anyway}, and what the pooled Bayesian likelihood
of \cref{sec:did:bayes} is forced to swallow into its residual $\sigma$. Hold on to
\eur{\nbEightAlphaSd}; it appears three more times, wearing three different costumes.

\emph{The effect is a clean step.} The app adds \eur{\nbEightTrueEffect} per store per week from
week \nbEightLaunch{} onward --- no build-up, no decay, no anticipation. That is the grade every
estimator in this chapter is marked against, and the last thing in the chapter that is known for
certain.

% =====================================================================================
\section{Identification}
\label{sec:did:ident}

\subsection*{The estimand}

In the potential-outcome notation of \citet{rubin1974} and \citet{imbens2015}, store $s$ in week
$t$ has two potential revenues: $Y_{st}(1)$ if the app is live in that store and $Y_{st}(0)$ if it
is not. Only one is ever observed. The object the chain is buying is the \emph{average treatment
effect on the treated} --- the effect in the stores that actually got the app, after they got it:
%
\begin{equation}
  \text{ATT}
  \;=\;
  \E\big[\, Y_{st}(1) - Y_{st}(0) \;\big|\; D_s = 1,\ t \ge t_0 \,\big].
  \label{eq:did:estimand}
\end{equation}
%
The final \textbf{T} is not a technicality and \cref{sec:did:euro} will spend real money on it.
\Cref{eq:did:estimand} is the effect \emph{on the pilot stores}. It is not the effect the other
\nbEightNOther{} stores of the chain would experience, and since pilot sites are chosen rather than randomised,
there is no reason for the two to coincide. Everything in this chapter estimates
\cref{eq:did:estimand}; nothing in this chapter estimates the chain-wide effect, and
\cref{sec:did:euro} treats that gap as the decision's dominant uncertainty rather than hiding it.

The second term inside \cref{eq:did:estimand} is the counterfactual: what the pilot stores would
have earned after week \nbEightLaunch{} had the app never shipped. Two assumptions license the
control stores to stand in for it.

\subsection*{The assumptions}

\begin{assumption}[Parallel trends]\label{as:did:pt}
Stated on \emph{untreated} potential outcomes: for any post-launch week $t \ge t_0$ and any
pre-launch week $t' < t_0$,
%
\begin{equation}
  \E\big[\, Y_{st}(0) - Y_{st'}(0) \;\big|\; D_s = 1 \,\big]
  \;=\;
  \E\big[\, Y_{st}(0) - Y_{st'}(0) \;\big|\; D_s = 0 \,\big].
  \label{eq:did:pt}
\end{equation}
%
Had the app never launched, the pilot stores' revenue would have \emph{changed} by the same amount
as the controls'. \Cref{eq:did:pt} says nothing whatever about \emph{levels}: pilot stores may be
systematically bigger, busier or richer, and DiD does not care, because a time-invariant level
difference cancels in the contrast. What must match is the change. In marketing terms: the seasonal
tide, chain-wide promotions and macro shocks hit both groups alike, and nothing \emph{time-varying
and group-specific} --- a regional campaign running only in pilot cities, a refit only in flagship
stores --- is in play.
\end{assumption}

\begin{assumption}[No anticipation]\label{as:did:na}
The app moves nothing before it exists:
%
\begin{equation}
  \E\big[\, Y_{st}(1) - Y_{st}(0) \;\big|\; D_s = 1 \,\big] \;=\; 0
  \qquad \text{for all } t < t_0 .
  \label{eq:did:na}
\end{equation}
%
This sounds vacuous and is not. Pre-launch buzz marketing, staff trained weeks early, or customers
\emph{deferring} purchases until the app's perks arrive all violate it --- and they contaminate the
very ``before'' baseline that \cref{eq:did:twobytwo} subtracts, so the estimator ends up
subtracting part of the app from itself. \Cref{sec:did:wrong} breaks \cref{as:did:na} on purpose and
measures the damage.
\end{assumption}

\begin{assumption}[No interference across stores]\label{as:did:sutva}
The app in a pilot store does not change revenue in any control store --- SUTVA
\citep{imbens2015, rubin1980}, applied to a store network. If app-holders defect from a nearby
control store, the control trend dips and DiD \emph{overstates} the lift; if pilot buzz lifts the
whole chain, it \emph{understates} it. No within-sample diagnostic can see either, and this chapter
does not pretend otherwise: \cref{as:did:sutva} is discharged by market-structure judgment (flag
pilot stores sharing a catchment with a control), not by a test.
\end{assumption}

\subsection*{Why the two differences give the ATT}

Granting \cref{as:did:na}, a treated store's pre-launch observation \emph{is} its untreated
potential outcome, so the treated group's observed change decomposes as
%
\begin{equation}
  \underbrace{\E\big[Y_{s,\text{post}} - Y_{s,\text{pre}} \mid D_s = 1\big]}_{\text{observed treated change}}
  \;=\;
  \text{ATT}
  \;+\;
  \underbrace{\E\big[Y_{s,\text{post}}(0) - Y_{s,\text{pre}}(0) \mid D_s = 1\big]}_{\text{the counterfactual change --- unobserved}} .
  \label{eq:did:decomp}
\end{equation}
%
\Cref{as:did:pt} says the unobserved second term equals the \emph{controls'} observed change.
Substituting it and rearranging gives exactly \cref{eq:did:twobytwo}. That is the whole derivation,
and it is worth seeing how little it contains: no functional form, no distributional assumption, no
likelihood. The causal content of this chapter is two lines of algebra and two assumptions. Every
estimator that follows --- classical or Bayesian --- is bookkeeping on top of them.

\begin{remark}[An assumption with no test is a wish]
\Cref{as:did:pt} is \textbf{untestable}. This has to be said in the plainest possible terms,
because the chapter's central diagnostic is routinely mistaken for a test of it.
\Cref{eq:did:pt} is a statement about the treated stores' \emph{unobserved} post-launch untreated
path. No data can reach it. What the event study of \cref{sec:did:classical} supplies is evidence
that the two groups \emph{did} move in parallel before launch, which is a different proposition, and
it is the only evidence anyone will ever have. A time-varying, group-specific shock arriving
exactly at week \nbEightLaunch{} would leave the pre-trend perfectly flat and destroy the estimate.
\Citet{roth2022} adds a sharper warning: conditioning the analysis on having \emph{passed} a
pre-trend test distorts the estimator that follows, because the panels that pass are exactly the
ones whose pre-period noise happened to run the convenient way. The correct posture is to report
the leads, report their power, and be honest that a flat pre-trend is a licence to proceed rather
than a proof of anything.
\end{remark}

% =====================================================================================
\section{The classical baseline}
\label{sec:did:classical}

Before any sampler, the discipline of this book is to ask what a competent analyst produces
\emph{without} one. Here the answer is not a different analysis: it is the same estimand
(\cref{eq:did:estimand}), identified by the same two assumptions, estimated with the simplest tool
that is still correct.

\subsection*{The estimator, naked}

The canonical $2 \times 2$ DiD is four cell means and two subtractions ---
\cref{eq:did:twobytwo}, computed directly. Run instead as a regression on the store-week panel,
%
\begin{equation}
  Y_{st}
  \;=\;
  \beta_0
  \;+\; \beta_1 D_s
  \;+\; \beta_2 \,\text{post}_t
  \;+\; \beta_3 \big(D_s \times \text{post}_t\big)
  \;+\; \varepsilon_{st},
  \label{eq:did:ols}
\end{equation}
%
the interaction coefficient $\beta_3$ \emph{is} that difference of differences, algebraically and
exactly --- no approximation, and the notebook asserts the identity to machine precision. Fitted to
the panel, it returns \eur{\nbEightAttCl} per store per week against a planted
\eur{\nbEightTrueEffect} --- \eur{\nbEightAttClErrAbs} low. The other two coefficients are worth
naming, because one of them is about to earn its keep. $\beta_2$ is the seasonal tide (the
controls' change). $\beta_1$ is the \textbf{level gap} --- how much more a pilot store sells than a
control store before anything happens. It is not the estimand, and in this simulated world, where
treatment was assigned at random, it is a near-nothing \eur{\nbEightLvlEst}. Keep it in view
anyway.

\subsection*{The covariance choice is the estimate's other half}

A point estimate without a defensible standard error is half an answer, and the \emph{default}
standard error here is indefensible. The \nbEightNObs{} rows are not \nbEightNObs{} independent
observations: they are \nbEightNStores{} stores watched for \nbEightNWeeks{} weeks, and a store's
weeks resemble each other --- they share the baseline $\alpha_s$ of \cref{eq:did:dgp}, and in a
real chain they would share serially correlated local shocks on top. The information content of the
panel is closer to \nbEightNStores{} than to \nbEightNObs.

Clustering on store is the classical fix \citep{liang1986, cameron2015}: it lets the errors be
arbitrarily correlated \emph{within} a store and asks only for independence \emph{across} stores,
so the \nbEightNStores{} stores, not the \nbEightNObs{} rows, become the unit of inference. That
correction is the point of \citet{bertrand2004}, whose famous finding --- that serial correlation
in panel data can leave DiD standard errors wrong by a factor of five or more --- reorganised the
applied literature.

Fit \cref{eq:did:ols} twice, once with iid (``nonrobust'') errors and once clustered on store. The
point estimate is identical to the last decimal; only the story about its precision changes. Read
\cref{tab:did:se} before reading the paragraph after it, because the direction of the change is not
the one the folklore prepares anyone for.

\input{build/tables/nb08_se}

\subsection*{The lesson is \emph{not} ``clustering widens intervals''}

On the estimand --- the interaction --- clustering \textbf{shrank} the standard error, from
\eur{\nbEightAttIidSe} to \eur{\nbEightAttClSe}, a factor of \nbEightSeRatioDid{}: the iid interval
was roughly twice too \emph{wide}. On the level gap from the \emph{identical} regression, clustering
\textbf{inflated} it, from \eur{\nbEightLvlIidSe} to \eur{\nbEightLvlClsSe}, a factor of
\nbEightSeRatioLevel: the iid interval was three times too \emph{narrow}. Both are correct, and the
mechanism is one object seen from two sides. The object is the store baseline $\alpha_s$, whose
standard deviation, \eur{\nbEightAlphaSd}, exceeds the weekly noise.

\begin{itemize}[leftmargin=*]
  \item The \textbf{level gap} compares \emph{different stores}, so it is fully exposed to the
        baseline scatter --- but the iid formula, believing it holds \nbEightNObs{} independent
        rows, divides that scatter by $\sqrt{\nbEightNObs}$ instead of by $\sqrt{\nbEightNStores}$.
        It reports a precision the data never had. This is the classic anti-conservative failure,
        and it is the one that ends careers: a confident interval around a number nobody can
        actually resolve.
  \item The \textbf{DiD interaction} compares each store \emph{against its own past}, so $\alpha_s$
        cancels out of the contrast entirely --- that is the bargain of \cref{sec:did:decision}.
        But it does \emph{not} cancel out of the pooled residual of \cref{eq:did:ols}, and iid OLS
        prices the estimate off that pooled residual: it charges the estimate for a scatter the
        contrast never used. The cluster-robust sandwich looks at the \nbEightNStores{} store-level
        trajectories directly, notices the cancellation, and hands back the honest --- here,
        \emph{smaller} --- standard error.
\end{itemize}

So the iid standard error is \textbf{not conservative and not anti-conservative as a rule. It is
wrong in a direction that cannot be predicted without doing the clustering, and the sign of the
error flips between two coefficients of the same fit.} That sentence is the entire argument for
making the covariance choice explicit rather than accepting a default, and it is a sharper argument
than the folk version of \citet{bertrand2004} that circulates in practice (``cluster, or your
$p$-values will be too small''). The folk version is a rule of thumb about a direction. The truth
is that there is no direction.

\subsection*{Bertrand--Duflo--Mullainathan's own case, isolated}

The DGP's weekly shocks are independent over time. Real store revenue is not: a good month tends to
be followed by another good month, and that persistence is precisely what \citet{bertrand2004}
studied. So rebuild the identical panel with AR(1) shocks of the same marginal standard deviation,
%
\begin{equation}
  \varepsilon_{st} \;=\; \rho \, \varepsilon_{s,t-1} \;+\; u_{st},
  \qquad
  u_{st} \sim \mathcal N\!\Big(0,\ 80^{2}\,(1 - \rho^{2})\Big),
  \label{eq:did:ar1}
\end{equation}
%
sweep the persistence $\rho$, and refit --- averaging over \nbEightAroneNPanels{} fresh panels per
$\rho$, since a single panel's standard error is itself noisy. Only the shocks change; the
estimator, the truth and the design do not. \Cref{tab:did:arone} reports the sweep.

\input{build/tables/nb08_arone}

Two readings, and the second is the one usually missed. First, the iid standard error moves by
\eur{\nbEightAroneIidSpread} across the entire sweep. It is not \emph{under}-estimating the
uncertainty that persistence creates so much as it is \emph{structurally unable to perceive it}:
serial correlation appears nowhere in the nonrobust formula. That blindness is the gap
\citeauthor{bertrand2004} measured, and it is why the estimator's true uncertainty can double while
its reported uncertainty sits unmoved.

Second, \textbf{the clustered column is not monotone in $\rho$}, and pretending otherwise would be
to teach the wrong mechanism. The honest standard error climbs from \eur{\nbEightAroneSeZero} at
$\rho = 0$ to \eur{\nbEightAroneSePeak} at $\rho = \nbEightAroneRhoPeak$ --- a factor of
\nbEightAroneSeGrowth{} --- and then \emph{eases} to \eur{\nbEightAroneSeTop} at
$\rho = \nbEightAroneRhoTop$. The reason is exactly the mechanism of the previous subsection: a
near-permanent shock behaves like a store baseline, and a store baseline \emph{differences out of
the DiD contrast}. It is \emph{intermediate} persistence --- shocks that outlive the pre/post
boundary but not the panel --- that does the most damage, because such a shock is neither fresh
noise the averaging can kill nor a fixed effect the differencing can cancel. Persistence is not a
dial that monotonically inflates uncertainty; it is a dial that inflates it and then, at the
extreme, hands it back.

\subsection*{The event study: the only evidence for parallel trends there will ever be}

The $2 \times 2$ collapses twelve post-launch weeks into one number. An \emph{event study} refuses
the collapse and estimates one coefficient per week relative to launch, from a two-way
fixed-effects regression with store and week effects,
%
\begin{equation}
  Y_{st}
  \;=\;
  \alpha_s \;+\; \lambda_t
  \;+\; \sum_{k \ne -1} \beta_k \, \mathbf 1[\,t - t_0 = k\,] \cdot D_s
  \;+\; \varepsilon_{st},
  \label{eq:did:eventstudy}
\end{equation}
%
with store-clustered standard errors. The $\beta_k$ for $k < 0$ are \emph{leads}; the $\beta_k$ for
$k \ge 0$ are \emph{lags}, and together they are the dynamic ATT. \Cref{fig:did:eventstudy} plots
them.

\input{build/floats/nb08_eventstudy}

Read the pre-period first, and read it carefully, because \cref{eq:did:eventstudy} contains a trap
that catches new readers every time. There are \nbEightEsNPre{} pre-launch weeks in this panel but
only \textbf{\nbEightEsNLeads{} leads}. The week $k = -1$ is spent as the \emph{reference category}:
every other coefficient is measured relative to it, so it is \textbf{pinned to zero by
construction, not by evidence}. A reader who counts twelve flat pre-launch dots and concludes that
twelve weeks of parallel trends have been verified has counted one dot that could not have been
anything else. The dot at $k = -1$ is a normalisation. It is not a finding, and it never will be.

The \nbEightEsNLeads{} leads that \emph{are} free to move do not move.
\nbEightEsLeadsOut{} of \nbEightEsNLeads{} lead intervals exclude zero, where 90\,\% intervals
would put roughly one there by chance alone; the mean absolute lead is
\eur{\nbEightEsMeanAbsLead} against a typical lead standard error of \eur{\nbEightEsLeadSe}. That
is noise, not a pre-trend. It is the honest content of the sentence ``parallel trends look
credible'', and it is worth being explicit that the evidence is a \emph{non-rejection}: the leads
are consistent with \cref{as:did:pt}, and they would also be consistent with a violation too small
for \nbEightNStores{} stores to resolve.

Then the effect steps. The \nbEightEsNLags{} lags average \eur{\nbEightEsLagMean} and run from
\eur{\nbEightEsLagMin} to \eur{\nbEightEsLagMax} --- an immediate, flat step, with no build-up and
no decay, which is what \cref{eq:did:dgp} planted. Averaging the lags reproduces the $2 \times 2$
interaction to within \eur{\nbEightEsLagVsDid}: the two estimators are the same estimator, read at
two resolutions.

\subsection*{What the classical interval does not say}

The clustered 90\,\% confidence interval is
$[\eur{\nbEightAttClLo},\ \eur{\nbEightAttClHi}]$, and the planted \eur{\nbEightTrueEffect} falls
inside it. It is \textbf{not} a 90\,\% probability that the true effect lies in that range. A
confidence interval is a property of the \emph{procedure} --- 90\,\% of intervals built this way
would cover the truth across repeated samples --- and this particular interval either covers the
truth or does not.

The distinction is not pedantry, and \cref{sec:did:euro} is exactly where it bites. The rollout
rule needs $\Prob(\text{ATT} > \eur{\nbEightAppCost})$: a probability \emph{about the effect}. The
classical apparatus cannot supply one --- not because its arithmetic is inferior, but because that
quantity does not exist in its vocabulary. It exists in a posterior. That, and not a better point
estimate, is what the next section is for.

% =====================================================================================
\section{The Bayesian model}
\label{sec:did:bayes}

\subsection*{One modelling choice, made for the right reason}

Before writing a likelihood, collapse each store's \nbEightNWeeks{} weekly observations into a
single pre-launch mean and a single post-launch mean --- \nbEightNCollapsed{} rows in place of
\nbEightNObs. This is the \citet{bertrand2004} guard, taken by a different route than
\cref{sec:did:classical}'s. Clustering \emph{prices} the within-store correlation; collapsing
\emph{removes} it, by refusing to treat \nbEightNWeeks{} weekly rows per store as \nbEightNWeeks{}
independent draws in the first place. Given \cref{tab:did:se}, where the sign of the iid error
flipped between two coefficients of one fit, removing the problem rather than betting on the
direction of its bias is the conservative move. It is a modelling decision, not an API constraint
--- though CausalPy's \code{DifferenceInDifferences} \citep{causalpy} does consume the classic
$2 \times 2$ form cleanly.

\subsection*{Likelihood and priors}

On design row $\mathbf x_i = (1,\ g_i,\ p_i,\ g_i p_i)^\top$ --- intercept, group flag, post flag,
interaction, with $i$ running over the \nbEightNCollapsed{} collapsed store-periods --- the model is
%
\begin{align}
  \tilde Y_i \mid \boldsymbol\beta, \sigma
    &\;\sim\; \mathcal N\!\big(\mathbf x_i^\top \boldsymbol\beta,\ \sigma^{2}\big),
    \label{eq:did:lik}\\
  \beta_j &\;\sim\; \mathcal N(0,\ 50^{2}), \label{eq:did:priorb}\\
  \sigma &\;\sim\; \text{HalfNormal}(1), \label{eq:did:priors}
\end{align}
%
and the ATT is $\beta_3$ --- the \emph{same} $\beta_3$ that \cref{eq:did:ols} estimated by ordinary
least squares. Same estimand, same identification, a different apparatus for expressing uncertainty
about it. That is what makes \cref{sec:did:compare} a comparison rather than a change of subject.

The priors \cref{eq:did:priorb,eq:did:priors} are \emph{fixed by the library} and scaled for data of
order one. On raw euro revenue --- levels around \eur{1000}, residual spread in the hundreds --- a
$\text{HalfNormal}(1)$ prior on $\sigma$ is absurdly tight: it would fight the likelihood and drag
the effect low. The remedy is standardisation, not a fight: revenue is centred and scaled before
the fit, and the effect is back-transformed to euros afterwards. This is worth naming because it is
the single most common way a Bayesian analysis silently goes wrong in practice --- a default prior
meeting data on a scale it was never written for. As throughout this book, the prior is doing very
little work once the scale is right; \cref{eq:did:lik} is doing all of it, and
\cref{sec:did:compare} is about precisely that.

Inference is by the No-U-Turn sampler \citep{hoffman2014} in PyMC \citep{salvatier2016}, with
$\hat R = \nbEightBayesRhat$, a minimum bulk effective sample size of \nbEightBayesEss{}
\citep{vehtari2021} and \nbEightBayesDiv{} divergent transitions. The chains have converged; nothing
that follows is a sampler artefact.

The posterior mean ATT is \eur{\nbEightBayesAtt} per store per week, with a 90\,\% credible interval
of $[\eur{\nbEightBayesLo},\ \eur{\nbEightBayesHi}]$. The point estimate agrees with the classical
one to within a few euros. The interval is \nbEightBayesOverCluster$\times$ wider than the clustered
classical interval, and that fact is not a curiosity. It is this chapter's finding, and
\cref{sec:did:valid,sec:did:compare} take it apart.

% =====================================================================================
\section{Diagnostics and validation}
\label{sec:did:valid}

\subsection*{Model criticism: can the model generate the data at all?}

Before reading an effect off \cref{eq:did:lik}, ask whether it could plausibly have produced the
panel. The model is deliberately coarse: four cell means and one shared Gaussian noise term. Every
store baseline $\alpha_s$ it does not describe has to live inside $\sigma$.
\Cref{fig:did:ppc} draws replicated datasets from the posterior predictive and overlays them on the
\nbEightNCollapsed{} observed store-period revenues. \pc{\nbEightPpcCov} of the observations fall
inside their 5--95\,\% replicate band, so there is no gross misfit.

\input{build/floats/nb08_ppc}

But notice \emph{how} the model passes. Its posterior residual standard deviation is
\eur{\nbEightBayesSigma} --- fat enough to swallow the between-store baseline scatter of
\eur{\nbEightAlphaSd} that four cell means cannot describe. The check passes, and the very thing
that lets it pass is what \cref{sec:did:compare} will indict. A posterior predictive check asks
whether the model can generate data like this; it does not ask whether the model's uncertainty about
a \emph{contrast} is honest. Those are different questions, and only the second one is about to cost
money.

\subsection*{Placebo in time}

\Cref{as:did:na}'s mirror image: run the identical $2 \times 2$ on the \emph{pre-period only},
splitting it into a fake before and after at week \nbEightPlaceboTimeWeek. The app did not exist
then, so the placebo effect must be approximately zero --- a large one would mean the design
manufactures effects out of seasonality. It comes to \eur{\nbEightPlaceboTime} with a standard error
of \eur{\nbEightPlaceboTimeSe}: within sampling noise of zero, against a real effect of
\eur{\nbEightAttCl}.

\subsection*{Placebo in group: a permutation null}

The time placebo asks whether the design manufactures effects out of \emph{seasonality}. Its
complement asks whether it manufactures them out of \emph{store-to-store noise}. Take only the
control stores --- the app never touched them, so the true effect there is exactly zero --- declare
a random half of them ``pseudo-treated'', and run the same $2 \times 2$ at the real launch week.
Repeating over \nbEightPermN{} random relabellings traces out the permutation null: the full
distribution of estimates this pipeline produces \emph{when nothing is going on}
(\cref{fig:did:permutation}).

\input{build/floats/nb08_permutation}

Two readings. The first is a $p$-value: store noise never produced anything larger than
\eur{\nbEightPermMax} in magnitude, while the real estimate sits \nbEightPermZ{} null standard
deviations out, for $p \le \nbEightPermP$ (the resolution floor of \nbEightPermN{} draws). The
second is more useful, and it is the one to file away: the null's spread, \eur{\nbEightPermSd}, is
an \textbf{assumption-light measurement of how far store noise alone can move this estimator} ---
no likelihood, no prior, no Gaussian anywhere. Noise moves it by about $\pm\eur{\nbEightPermTwoSd}$ at two standard
deviations. The pooled Bayesian interval of \cref{sec:did:bayes} is $\pm\eur{\nbEightBayesHalf}$.
Something is wrong, and it is not the classical arm.

\subsection*{Recovery and calibration across fresh panels}

A single panel is a single draw of the world. The honest questions about an estimator are whether it
is \emph{centred} on the truth across repeated samples and whether its nominal 90\,\% interval
\emph{covers} at 90\,\%. Re-running the whole procedure on \nbEightSeedN{} fresh panels from
\cref{eq:did:dgp} answers both.

The point estimate is fine: mean \eur{\nbEightSeedMean} against a planted \eur{\nbEightTrueEffect},
a bias of \eur{\nbEightSeedBiasAbs}. The interval is not. It covers the truth on
\nbEightSeedCov{} of \nbEightSeedN{} panels --- which sounds like a triumph until the width is
compared with the scatter it is supposed to describe. The seed-to-seed standard deviation of the
estimate is \eur{\nbEightSeedSd}. The interval's half-width is \eur{\nbEightBayesHalf}, or
\textbf{\nbEightPooledOverScatter$\times$ the true sampling scatter}, where a calibrated 90\,\%
interval would sit at roughly \nbEightCalibratedRatio$\times$. This is \emph{over}-coverage: an
interval that is honest and enormously conservative. Coverage can fail in both directions, and
counting hits alone would have missed it entirely.

\subsection*{The diagnosis, and the repair}

The cause is visible in \cref{eq:did:lik} and was already named by \cref{fig:did:ppc}: the pooled
$2 \times 2$ model puts the store baselines $\alpha_s$ into the residual $\sigma$, and then charges
the effect for a scatter the DiD contrast never touches. The cure is the oldest trick in the
fixed-effects book --- difference each store against itself. Collapse every store to a single
number,
%
\begin{equation}
  \Delta_s \;=\; \bar Y_{s,\text{post}} - \bar Y_{s,\text{pre}},
  \label{eq:did:delta}
\end{equation}
%
which removes $\alpha_s$ \emph{exactly} (it is time-invariant, so it appears in both means and
cancels) while keeping the two-period \citeauthor{bertrand2004} guard. The DiD becomes a
one-coefficient comparison of store-level \emph{changes}:
%
\begin{equation}
  \Delta_s \;\sim\; \mathcal N\!\big(\gamma_0 + \gamma_1 D_s,\ \sigma_\Delta^{2}\big),
  \qquad
  \widehat{\text{DiD}} \;=\; \gamma_1 \;=\; \bar\Delta_{D=1} - \bar\Delta_{D=0},
  \label{eq:did:deltamodel}
\end{equation}
%
with $\gamma_0$ the control stores' mean change (the seasonal tide) and $\sigma_\Delta$ containing
only averaged week-level noise --- no store baselines. Same estimand, same point estimate, a
radically smaller residual.

It works, and --- the part that makes the tightening legitimate --- it works without breaking
coverage. The within-store posterior returns \eur{\nbEightFeAtt}, 90\,\% interval
$[\eur{\nbEightFeLo},\ \eur{\nbEightFeHi}]$, half-width \eur{\nbEightFeHalf}: about
\nbEightFeTighter$\times$ tighter than the pooled model's ($\hat R = \nbEightFeRhat$, minimum ESS
\nbEightFeEss, \nbEightFeDiv{} divergences). Re-run on the same \nbEightSeedN{} panels, it covers on
\nbEightFeCov{} of them, its seed-to-seed scatter is \eur{\nbEightFeSeedSd}, and its half-width is
therefore \nbEightFeHwOverScatter$\times$ that scatter --- against the calibrated benchmark of
\nbEightCalibratedRatio$\times$, and against the pooled model's \nbEightPooledOverScatter$\times$.
The two ratios are like for like: each divides a headline 90\,\% half-width by the sampling
standard deviation of its own estimator across the same \nbEightSeedN{} panels. (The refits'
\emph{mean} half-width, \eur{\nbEightFeSeedHalf}, sits just below the headline panel's, so the
comparison is not an artefact of the one panel that was drawn.)
\textbf{A tighter interval is progress only if coverage survives};
here it does, and the width that was removed was redundancy, not protection.

Note what has just happened. Three independent witnesses --- the permutation null's spread
(\eur{\nbEightPermSd}), the seed-to-seed scatter (\eur{\nbEightSeedSd}) and the within-store
posterior's half-width (\eur{\nbEightFeHalf}) --- now agree on the honest scale of this estimator's
uncertainty. And a fourth witness agreed with all of them from outside the Bayesian apparatus
entirely, before any of the others were computed: \cref{sec:did:classical}'s cluster-robust standard
error, which put the half-width at \eur{\nbEightAttClHalf} without a single prior.

% =====================================================================================
\section{Point estimate versus posterior}
\label{sec:did:compare}

\input{build/tables/nb08_compare}
\input{build/floats/nb08_forest}

\subsection*{On the number, they agree}

\Cref{tab:did:compare} and \cref{fig:did:forest} put five estimators on the estimand of
\cref{eq:did:estimand}: two classical, two Bayesian, and the deliberately-wrong iid row. The point
estimates span \eur{\nbEightEstSpan} --- \pc{\nbEightEstSpanPct} of the effect itself --- and every
interval covers the planted truth.

That is not a coincidence and it is not a disappointment. \textbf{The causal work was done by the
identification, not by the machinery.} Parallel trends plus no anticipation is what turns a
difference of averages into an effect; once that argument is granted, a Bayesian regression with
weak priors and \nbEightNObs{} observations is doing arithmetic that four cell means had already
done. This is the Bernstein--von Mises phenomenon \citep{gelman2013}, and it is the ordinary case in
this book: with a well-scaled likelihood and plentiful data, the posterior concentrates where the
maximum-likelihood estimate does. Anyone hoping that ``going Bayesian'' would move the point
estimate should read this chapter as the correction.

\subsection*{On the width, they do not --- and here Bayes lost}

The classical clustered interval (half-width \eur{\nbEightAttClHalf}) and the within-store posterior
(\eur{\nbEightFeHalf}) land on the same scale. The pooled Bayesian $2 \times 2$ is the outlier:
half-width \eur{\nbEightBayesHalf}, about \nbEightBayesOverCluster$\times$ wider than either.

State it bluntly, because it is load-bearing for the whole book. \textbf{The
competently-executed classical estimator had the right interval on the first attempt; it took until
\cref{eq:did:deltamodel} --- a second model, written after a multi-seed calibration study had
diagnosed the first one --- to get the Bayesian arm back to the same place.} The classical route
forced the covariance assumption to be stated out loud, and therefore forced it to be got right. The
Bayesian route let a coarse likelihood swallow the store baselines into $\sigma$ and then report the
consequences with a perfectly straight face, converged chains and a clean posterior predictive check.
A posterior inherits every modelling sin, silently.

\subsection*{Misspecification goes both ways --- which is why it is specification, not paradigm}

Now put this chapter next to \refgeo, because together they make the argument neither could
make alone.

In \refgeo, a Bayesian synthetic control applied an iid likelihood to errors that were in fact
a persistent random walk. It under-priced the standard deviation of a twenty-week sum by a factor of
about \nbSevenUnderprice, and its nominal 90\,\% interval covered the truth on
\pc{\nbSevenTotalCov} of fresh panels. The posterior was \textbf{far too narrow}.

In this chapter, a Bayesian $2 \times 2$ applied a pooled likelihood to a contrast that had already
differenced the store baselines away. Its interval came out \nbEightBayesOverCluster$\times$ wider
than the clustered classical one --- interval against interval, the like-for-like comparison ---
and \nbEightPooledOverScatter$\times$ the true sampling scatter where a calibrated interval belongs
at \nbEightCalibratedRatio$\times$; it covered on \nbEightSeedCov{} of \nbEightSeedN{} panels. The
posterior was \textbf{far too wide}.

Same paradigm. Same weak priors. Same converged chains. Opposite errors. If Bayesian inference were
inherently overconfident, \cref{ch:did} could not have happened; if it were inherently conservative,
\refgeo{} could not have. \textbf{What failed in each case was a specific, nameable, repairable
modelling choice about the error structure --- and the sign of the resulting error is a property of
that choice, not of the paradigm.} The prior did almost nothing in either chapter. The likelihood was
the whole ballgame in both.

The same reading explains why the classical robustness toolkit won both times, and it is not because
frequentism is deeper. Cluster-robust and HAC standard errors, the bootstrap and randomization
inference all \emph{decline to assume a shape for the errors} and go and measure one instead
\citep{liang1986, newey1987, bertrand2004, cameron2015}. The cluster-robust sandwich of
\cref{sec:did:classical} did not need to be told that the store baselines cancel in the DiD contrast
and do not cancel in the level gap; it looked at the \nbEightNStores{} store trajectories and found
out. \textbf{It won because it committed to less.} And where the Bayesian scale is the one that must
be trusted, the principled repair is to correct it rather than abandon it: \citet{mueller2013} shows
a posterior computed under a misspecified likelihood can be rescaled by a sandwich covariance so that
its intervals recover their frequentist coverage --- the centre is fine, only the scale is wrong,
in both chapters and in both directions.

\subsection*{What the posterior did buy}

One thing, and it is decisive. Look at the last column of \cref{tab:did:compare}. The classical rows
read \emph{not defined}, and that is the literal truth rather than a rhetorical concession: the
frequentist apparatus has no probability to attach to a hypothesis about a fixed parameter. Yet the
rollout rule of \cref{sec:did:euro} is nothing but such a probability. So the decision quantity ---
$\Prob(\text{ATT} > c)$ --- lives in exactly one of the two columns, and the euro decision consumes
nothing else.

The ledger, then, stated plainly. Bayes did not buy the point estimate: the classical one is at
least as good. It did not buy the interval: on the first attempt it bought a much worse one, and it
took a second model to recover the classical arm's width. What it bought is a coherent way to
propagate uncertainty into euros, and the only object in the chapter that can answer the question
the chain actually asked. \textbf{Use the posterior for the decision, and design-based inference for
the confidence in it.}

% =====================================================================================
\section{The decision, in euros}
\label{sec:did:euro}

The vendor wants \eur{\nbEightAppCost} per store per week. The rule, a stated convention applied to
the posterior:
%
\begin{equation}
  \Prob(\text{ATT} > c)\ \begin{cases}
    > 0.9 & \Rightarrow\ \textsc{roll out},\\
    < 0.5 & \Rightarrow\ \textsc{no-go},\\
    \text{otherwise} & \Rightarrow\ \textsc{pilot further}.
  \end{cases}
  \label{eq:did:rule}
\end{equation}
%
At $c = \eur{\nbEightAppCost}$, and deliberately using the \emph{conservative} pooled posterior of
\cref{sec:did:bayes} rather than the tight within-store one --- a rollout case that survives the wide
interval is a robust case --- the arithmetic is emphatic. Net value is \eur{\nbEightNetMean} per
store per week; across \nbEightNChain{} stores and \nbEightWeeksYear{} weeks that is
\eur{\nbEightAnnualMeanM}\,m a year, with a 90\,\% interval running from
\eur{\nbEightAnnualLoM}\,m to \eur{\nbEightAnnualHiM}\,m. And
$\Prob(\text{the app pays}) = \nbEightPPays$. The verdict is \textbf{\nbEightDecision}
(\cref{fig:did:euro}).

\input{build/floats/nb08_euro}

\subsection*{The number that is not on the slide}

And that is where a careless analysis stops, having answered the wrong question with great
precision. \Cref{eq:did:estimand} is the ATT: the effect \textbf{on the pilot stores}. Rolling the
app out to all \nbEightNChain{} stores assumes the other stores respond the same way, and chains
rarely guarantee that, because pilot sites are not picked at random. The pilot contains \emph{no
information whatever} about the transfer share, and no amount of statistical sophistication applied
to the pilot will create any.

So the estimate is stressed on two axes at once: the running cost (which the vendor sets) and the
transfer share (which nobody has measured). \Cref{tab:did:stress} is the result, and it splits
cleanly.

\input{build/tables/nb08_stress}

The first thing the table says is that \emph{at the price actually on the table nothing collapses}.
Read down the \eur{\nbEightAppCost} column: the call survives every transfer assumption in it, and
even a world in which the other \nbEightNOther{} stores realise only half the pilot lift still
returns $\Prob(\text{pays}) = \nbEightPHalfAppCost$. The vendor's quote is not close to the
boundary, and no reading of the table makes it close.

What the table does say is that the \emph{two axes interact}, and neither alone decides. At full
transfer the cost may climb a long way before the call moves --- at \eur{\nbEightCostHigh} per
store-week it is still $\Prob(\text{pays}) = \nbEightPFullCostHigh$. But the room to move shrinks
with the transfer share: at half transfer, \eur{\nbEightCostMid} is already a coin flip
($\nbEightPHalfCostMid$) and \eur{\nbEightCostHigh} is a flat no ($\nbEightPHalfCostHigh$).
\textbf{The transfer share does not decide the rollout; it decides how much the chain can afford to
pay for it} --- and the transfer share is the one quantity the pilot contains no information about.

The honest deliverable is therefore not the base case's \eur{\nbEightAnnualMeanM}\,m. It is the
break-even band: the running cost the rollout still clears at 90\,\% posterior confidence, quoted at
each transfer assumption. Priced off the pooled posterior of \cref{sec:did:bayes} it runs from
\eur{\nbEightCapFull} per store-week at full transfer, to \eur{\nbEightCapThreeQuarter} at
three-quarters, to \eur{\nbEightCapHalf} at half. Hold the vendor's quote against that band --- with
the correction the next subsection prices. If the quote lands inside it, the cheap de-risking move
is a second pilot wave in deliberately \emph{non}-pilot-like stores, which measures the transfer
share directly instead of assuming it.

\subsection*{A note on which posterior priced this}

That band was priced off the posterior this chapter has just spent a section condemning, and the
conservatism has a price of its own. \Cref{sec:did:valid} found the pooled interval
\nbEightBayesOverCluster$\times$ wider than the clustered classical one; a ceiling read off its 10th
percentile inherits the whole of that excess width. Priced off the \emph{calibrated} within-store
posterior of \cref{eq:did:deltamodel} --- same point estimate, honest width, coverage verified over
\nbEightSeedN{} panels --- the band runs from \eur{\nbEightCapFullFe} at full transfer to
\eur{\nbEightCapThreeQuarterFe} at three-quarters and \eur{\nbEightCapHalfFe} at half.

The gap is the cost of the conservatism, and it is a decision, not a rounding. A quote landing
between \eur{\nbEightCapFull} and \eur{\nbEightCapFullFe} per store-week is a \textsc{no-go} on the
pooled band and a \textsc{roll out} on the calibrated one, on identical data: what changed between
those two verdicts is not the world but the residual the model was willing to charge the effect for.
The same fork appears in the break-even cost --- \eur{\nbEightBreakEven} pooled against
\eur{\nbEightBreakEvenFe} calibrated --- and in the probabilities near the effect itself: at
\eur{\nbEightCostNearAtt}, $\Prob(\text{ATT} > \eur{\nbEightCostNearAtt})$ is \nbEightPPooledNearAtt{}
on the pooled model and \nbEightPFeNearAtt{} on the calibrated one. \textbf{Far from the decision
boundary a conservative interval costs nothing; near it, a conservative interval is a decision.}
Quote the calibrated band, and keep the pooled one only as the stress case. That is the practical
reason \cref{sec:did:valid}'s calibration work is not an academic exercise.

% =====================================================================================
\section{What can go wrong}
\label{sec:did:wrong}

\subsection*{Parallel trends}

\Cref{as:did:pt}, broken on purpose --- and broken in the way a real chain would break it. The pilot
sites were the flagships, and flagships were already pulling away from the rest of the estate before
anyone wrote an app. The simulator plants exactly that: the same \nbEightNStores{} stores, the same
launch, the same true \eur{\nbEightTrueEffect}, plus a linear trend of \eur{\nbEightPtSlope} per week
carried by the pilot stores alone. Nothing else moves. Levels are still free (DiD never cared about
them); what is now false is that the two groups would have \emph{changed} alike.

\input{build/floats/nb08_pretrend}

The event study --- \cref{eq:did:eventstudy}, the same estimator with the same $k = -1$ reference
week and the same store-clustered errors as \cref{fig:did:eventstudy} --- catches it, and catches it
in the one way that cannot be mistaken for noise: the leads do not scatter, they \emph{slope}.
\Cref{fig:did:pretrend} fits \eur{\nbEightPtSlopeHat} per week through them against a planted
\eur{\nbEightPtSlope}, and \nbEightPtLeadsOut{} of the \nbEightPtNLeads{} free leads exclude zero
where roughly one would by chance. A flat pre-period is a scatter about zero; a violated one has a
direction.

The damage is arithmetic. The pre-window and the post-window are \nbEightLaunch{} weeks apart in
mean event time, so \nbEightLaunch{} weeks of divergence at \eur{\nbEightPtSlope} per week are booked
as app effect: the naive $2 \times 2$ reports \eur{\nbEightPtDid} against the true
\eur{\nbEightTrueEffect}, an \emph{upward} bias of \eur{\nbEightPtBiasAbs}. At the vendor's
\eur{\nbEightAppCost} the sign of the decision survives, which is precisely what makes the violation
dangerous: it does not announce itself by producing an absurd number.

The repair is to widen the model rather than the interval. Letting the two groups carry
\emph{different} linear trends --- a single group-by-week regressor added to
\cref{eq:did:eventstudy}, identified off the \nbEightEsNPre{} pre-launch weeks --- recovers
\eur{\nbEightPtFixed}. But the repair should be read for what it costs as well as what it buys. It
does not restore \cref{as:did:pt}; it \emph{replaces} it with a strictly stronger and equally
untestable one --- parallel trends conditional on a linear differential trend --- and a pre-trend
that is not linear will not be removed by a linear correction \citep{roth2022}. The
honest posture is the one \cref{sec:did:ident}'s remark already argued for: the leads are evidence,
not proof, and a design that needs a trend correction to work is a design whose control group was
chosen badly.

\subsection*{Anticipation}

\Cref{as:did:na}, broken on purpose. Suppose the pilot stores start behaving differently
\emph{before} week \nbEightLaunch{} --- pre-launch promo buzz, staff trained early, customers
holding back purchases until the app's perks arrive --- so that the lift effectively begins
\nbEightAntWeeks{} weeks early. The ``before'' baseline is then contaminated with effect, and
\cref{eq:did:twobytwo} subtracts part of the app from itself.

\input{build/floats/nb08_anticipation}

\Cref{fig:did:anticipation} catches it, and the panel is deliberately drawn a different way from
\cref{fig:did:eventstudy,fig:did:pretrend}. It plots the raw per-week treated-minus-control gap,
centred on the pre-period mean, with \textbf{no reference category at all} --- every one of the
\nbEightAntNPre{} pre-launch weeks is free to move. That choice is the point. Under anticipation the
week $k = -1$ is itself contaminated, so normalising on it would bury part of the violation inside
the normalisation and quietly rescale the rest; a diagnostic for a broken baseline must not be
anchored to the broken baseline. The price of the choice is that the dots no longer average to the
regression's $\beta_k$ --- they average to zero by construction --- and their bands are
cross-sectional rather than clustered, so they are wider than \cref{fig:did:eventstudy}'s.

The weeks $k = -2$ and $k = -1$ jump to \eur{\nbEightAntLeadMinusTwo} and
\eur{\nbEightAntLeadMinusOne}, far outside their bands, and because centred gaps must average zero
the contamination drags the \emph{other} pre-launch weeks negative too: the tell is a broken, tilted
pre-period, not merely two high dots. \nbEightAntLeadsOut{} of \nbEightAntNPre{} pre-launch weeks now
exclude zero, against roughly one expected by chance. The naive $2 \times 2$ reports
\eur{\nbEightAntDid} against a true \eur{\nbEightTrueEffect} --- a \emph{downward} bias of
\eur{\nbEightAntBiasAbs}, which is about $\nbEightAntWeeks/\nbEightLaunch$ of the effect: exactly the
share of the pre-period the contamination ate.

The pre-period check is therefore no formality: it just caught a violation that quietly removed a
sixth of the measured effect. The remedies are cheap once the violation is \emph{seen}: drop the
anticipation window from the pre-period (which restores \eur{\nbEightAntFixed}), or date the
treatment at \emph{announcement} rather than at launch. What the pre-period still cannot do is
certify \cref{as:did:pt}; a time-varying, group-specific shock arriving exactly at $t_0$ --- unlike
the sloping trend above, which the leads did catch --- remains invisible.

\subsection*{The two-way fixed-effects trap}

Real chains rarely launch everywhere at once, and this is where the chapter's title earns itself.
With \emph{staggered} adoption --- store $s$ adopting at its own week $g_s$ --- the workhorse
regression becomes static two-way fixed effects,
%
\begin{equation}
  Y_{st} \;=\; \alpha_s \;+\; \lambda_t \;+\; \beta^{\text{TWFE}} D_{st} \;+\; \varepsilon_{st},
  \qquad D_{st} = \mathbf 1[\,t \ge g_s\,],
  \label{eq:did:twfe}
\end{equation}
%
and \cref{eq:did:twfe} is a trap. \Citet{goodmanbacon2021} proved that
$\hat\beta^{\text{TWFE}}$ is a variance-weighted average of \emph{every} pairwise $2 \times 2$ DiD
the panel contains --- early-versus-late, treated-versus-never-treated, and, fatally,
\emph{late-versus-early}: comparisons that use \textbf{already-treated} cohorts as controls. The
Bacon weights themselves are variance shares and are non-negative; the damage arrives through what
those forbidden comparisons \emph{estimate}. When effects grow with time since adoption, a
comparison that uses an already-treated cohort as its control subtracts that cohort's still-growing
lift, and the estimate can land far from every true effect in the panel even though every one of them
is positive. Stated on the underlying cells rather than on the $2 \times 2$s,
\citet{dechaisemartin2020} show the same thing as a sign: TWFE assigns \textbf{negative implicit
weights} to some cohort-week effects, and here it assigns them to the largest ones
\citep{sunabraham2021}.

The notebook plants exactly that world: \nbEightStagNStores{} stores in two cohorts adopting at
weeks \nbEightStagCohortA{} and \nbEightStagCohortB, with a true effect that starts at
\eur{\nbEightStagBase} and \emph{grows} by \eur{\nbEightStagGrowth} per week after adoption.
\Cref{tab:did:staggered} reports what happens.

\input{build/tables/nb08_staggered}
\input{build/tables/nb08_bacon}

A single post dummy returns \eur{\nbEightTwfeNaive} against a true average post-adoption effect of
\eur{\nbEightTwfeTrue} --- \eur{\nbEightTwfeBiasAbs} low, and, worse than merely biased, it lands
\emph{below every single true post-adoption effect in the panel}, which start at
\eur{\nbEightStagBase} and rise from there. The mechanism is worth spelling out. Cohort A adopts at
week \nbEightStagCohortA{} and its effect keeps growing. When cohort B adopts at week
\nbEightStagCohortB, part of \cref{eq:did:twfe}'s estimate compares B against already-treated A ---
a ``control'' whose revenue is still rising with A's growing effect. B's lift is measured against
that rising counterfactual and shrunk, while the growing tail of A's effect looks like a calendar
trend that the week effects $\lambda_t$ partially absorb.

\Cref{tab:did:bacon} takes the coefficient apart rather than asserting the mechanism. With two
adoption dates and no never-treated stores, \cref{eq:did:twfe} has exactly two $2 \times 2$
comparisons to average, so the decomposition can be enumerated in full --- and the certificate that
it \emph{is} full is an identity: the weighted contributions reproduce the TWFE coefficient exactly,
\eur{\nbEightBaconSum} against \eur{\nbEightTwfeNaive}, to numerical precision. Read the two rows.
The clean comparison --- cohort A against not-yet-treated B --- returns \eur{\nbEightBaconDidClean}
and carries weight \nbEightBaconWClean; the forbidden one --- cohort B against already-treated A ---
returns \eur{\nbEightBaconDidForbidden} and carries weight \nbEightBaconWForbidden.

Two things follow, and the second is the one the folklore garbles. First, \textbf{both weights are
positive}, and the larger of the two sits on the forbidden comparison: what drags the coefficient
down is not a negative weight on that row but its \emph{estimate}, which measures B against a
control whose own effect is still climbing. Second, the negativity is real but it lives one level
down, on the cells rather than on the $2 \times 2$s: the implicit weight \cref{eq:did:twfe} places on
each treated cohort-week is negative for \nbEightBaconNegCells{} of the \nbEightBaconTreatedCells{}
treated cells, and they are precisely the cells carrying the panel's largest true effects. That is
\citeauthor{dechaisemartin2020}'s sign, computed rather than quoted, and it is why the coefficient
lands below every effect it is supposed to be averaging.

\textbf{Nothing in the TWFE output announces any of this.} The regression converges, the
coefficient has a standard error --- \eur{\nbEightTwfeSe}, from the store bootstrap of
\cref{tab:did:staggered}, comfortably tight enough to be reported with confidence --- and the $R^2$
is respectable. The bias is \eur{\nbEightTwfeBiasAbs}, or \nbEightTwfeBiasInSe{} bootstrap standard
errors: it is arithmetic, not sampling noise, and no amount of extra data would shrink it. That is
why the fix is a discipline rather than a diagnostic.

\subsection*{Callaway--Sant'Anna: never use the already-treated as controls}

\Citet{callaway2021} turn the diagnosis into a procedure. Estimate one clean $2 \times 2$ per cohort
$g$ and week $t$, anchored at the cohort's last untreated week $g - 1$, using \emph{only} stores not
yet treated at $t$ as controls:
%
\begin{equation}
  \widehat{\text{ATT}}(g, t)
  \;=\;
  \big[\bar Y_{g,t} - \bar Y_{g,\,g-1}\big]
  \;-\;
  \big[\bar Y_{\mathcal C(t),\,t} - \bar Y_{\mathcal C(t),\,g-1}\big],
  \qquad
  \mathcal C(t) = \{\, s : g_s > t \,\},
  \label{eq:did:cs}
\end{equation}
%
then aggregate the $(g,t)$ cells into whatever summary the decision needs --- an overall ATT, an
event-study profile, a per-cohort path. \Citet{sunabraham2021} reach the same destination by a
different route, interacting cohort with event time so that the event-study coefficients cannot
contaminate each other. On the planted panel, \cref{eq:did:cs} recovers \eur{\nbEightCsAtt} against a
true \eur{\nbEightCsTrue} --- within \eur{\nbEightCsErrAbs} of it --- from \nbEightCsCells{} clean
cells (\cref{fig:did:staggered}), and its
cohort-time effects track the growing truth rather than averaging it into a single misleading
scalar.

That \eur{\nbEightCsErrAbs} deserves a standard error rather than applause, and
\cref{tab:did:staggered} carries one: a cluster bootstrap over the stores, resampling
\emph{stores} rather than store-weeks, which is \citeauthor{callaway2021}'s own inference in its
simplest honest form. It puts the Callaway--Sant'Anna estimate at
\eur{\nbEightCsAtt} $\pm$ \eur{\nbEightCsSe}, a 90\,\% interval of
$[\eur{\nbEightCsLo},\ \eur{\nbEightCsHi}]$ that contains the truth --- so the miss is one panel's
noise, as it should be. The naive TWFE coefficient's bootstrap standard error is
\eur{\nbEightTwfeSe}, which is the whole point: its error is not a wide interval but a wrong centre,
and the interval is narrow enough to state the wrong centre with conviction.

\input{build/floats/nb08_cohorts}

One \textbf{design lesson} falls straight out of \cref{eq:did:cs}, and it is the most actionable
sentence in this chapter. The control set $\mathcal C(t)$ is \emph{empty} once every store has
adopted: after the last cohort switches on, there are no clean controls left and
$\widehat{\text{ATT}}(g,t)$ simply does not exist. \textbf{A real rollout must keep a never-treated
holdout} --- and it must keep it to the very end. \Cref{tab:did:staggered} prices the half-measure,
and prices it as a controlled contrast: the holdout panel is the \emph{same} \nbEightStagNStores{}
stores, drawn from the same seed in the same order, with \nbEightStagNeverTreated{} never-treated
stores appended, so the only thing that differs between the two TWFE rows is the presence of the
holdout. Adding it improves naive TWFE from \eur{\nbEightTwfeNaive} to \eur{\nbEightTwfeHoldout}
($\pm$ \eur{\nbEightTwfeHoldoutSe}), but it does \emph{not} fix it, because the forbidden
late-versus-early comparisons --- cohort B measured against already-treated A --- are still inside
the average. A holdout is necessary and not sufficient; the holdout buys the \emph{ability} to
estimate, and \cref{eq:did:cs} is what turns that ability into a number. An all-at-once rollout
destroys both, permanently, for the price of a few weeks' delay it was trying to save.

\subsection*{Spillover, and the limits}

\Cref{as:did:sutva} remains untested and untestable from this panel, and the chapter says so rather
than dressing a diagnostic in its clothes. It is a market-structure judgment: flag pilot stores that
share a catchment with a control, and either drop them or drop their neighbours before fitting
anything.

Two further limits are properties of the design rather than defects of the code. \Cref{as:did:pt} is
untestable and remains so however flat the leads are; \citet{roth2022} shows that even the act of
\emph{pre-testing} it distorts what follows. And the ATT is not the chain-wide effect, which
\cref{tab:did:stress} has already made the decision's dominant uncertainty. Both are reasons to lead
the conclusions with the event study and the transfer band rather than with a single euro figure.

% =====================================================================================
\section{Takeaways}
\label{sec:did:take}

\begin{enumerate}[leftmargin=*, itemsep=.45em]
  \item \textbf{The confounder may stay unobserved --- if it is constant.} DiD compares each store
        against its own past, so any time-invariant store characteristic differences out of
        \cref{eq:did:twobytwo} exactly. That is the bargain, and its price is \cref{as:did:pt},
        which is untestable.

  \item \textbf{The event study is the only evidence for parallel trends there will ever be --- and
        one of its dots is not evidence at all.} There are \nbEightEsNPre{} pre-launch weeks and
        \nbEightEsNLeads{} leads: $k = -1$ is the reference category, \emph{pinned to zero by
        construction}. Of the \nbEightEsNLeads{} free leads, \nbEightEsLeadsOut{} exclude zero
        (\cref{fig:did:eventstudy}). Report that as a non-rejection, not as a proof, and remember
        \citeauthor{roth2022}'s warning that pre-testing itself distorts what follows.

  \item \textbf{A violated pre-trend slopes; it does not scatter.} Give the pilot stores a trend of
        their own --- \eur{\nbEightPtSlope} a week, the flagship story --- and the leads pick it up
        at \eur{\nbEightPtSlopeHat} a week, with \nbEightPtLeadsOut{} of \nbEightPtNLeads{} excluding
        zero (\cref{fig:did:pretrend}), while the naive $2 \times 2$ books \nbEightLaunch{} weeks of
        pre-existing divergence as effect and reports \eur{\nbEightPtDid} against a true
        \eur{\nbEightTrueEffect}. Group-specific linear trends recover \eur{\nbEightPtFixed}, but
        they buy it by \emph{swapping} \cref{as:did:pt} for a stronger, equally untestable one.

  \item \textbf{The iid standard error is not conservative. It is wrong in an unpredictable
        direction.} On the DiD interaction, clustering \emph{shrank} the standard error by
        \nbEightSeShrinkDid$\times$ --- the store baseline cancels from the contrast but not from
        the pooled residual, so iid OLS charges the estimate for scatter it never used. On the level
        gap of the \emph{same} regression, clustering \emph{inflated} it \nbEightSeRatioLevel$\times$
        --- that contrast compares different stores, and iid divides the baseline scatter by
        $\sqrt{\nbEightNObs}$ instead of $\sqrt{\nbEightNStores}$ (\cref{tab:did:se}). The sign of
        the error flips inside one fit. That is the argument for making the covariance choice
        explicit.

  \item \textbf{Serial correlation does the most damage at \emph{intermediate} persistence.} The
        honest standard error climbs from \eur{\nbEightAroneSeZero} at $\rho = 0$ to
        \eur{\nbEightAroneSePeak} at $\rho = \nbEightAroneRhoPeak$ and then \emph{eases} at
        $\rho = \nbEightAroneRhoTop$, because a near-permanent shock behaves like a store baseline
        and differences out (\cref{tab:did:arone}). The iid standard error moves by
        \eur{\nbEightAroneIidSpread} across the whole sweep: it cannot perceive persistence at all.

  \item \textbf{Where Bayes lost --- and it lost badly, on the first attempt.} The pooled
        $2 \times 2$ posterior's interval was \nbEightBayesOverCluster$\times$ wider than either
        classical arm and \nbEightPooledOverScatter$\times$ the true seed-to-seed scatter, because
        the store baselines were left in $\sigma$. The competently-executed classical estimator had
        the right interval on the first attempt; it took a second Bayesian model
        (\cref{eq:did:deltamodel}, which differences $\alpha_s$ out) to get back to the same place.
        A posterior inherits every modelling sin, silently, with converged chains and a clean PPC.

  \item \textbf{Misspecification goes both ways --- which proves it is specification, not
        paradigm.} \Refgeo{}'s Bayesian interval was too \emph{narrow} --- it under-priced the
        standard deviation of a twenty-week sum by \nbSevenUnderprice$\times$; this chapter's was
        \nbEightBayesOverCluster$\times$ too \emph{wide}, interval against the clustered classical
        interval. Same paradigm, same weak priors, opposite errors. The prior did almost nothing in either chapter;
        the likelihood was the whole ballgame in both. The repair is to fix the likelihood
        (\cref{eq:did:deltamodel}) or to rescale the posterior \citep{mueller2013} --- not to lose
        faith in Bayes. And the classical toolkit won both times because it \emph{committed to
        less}: it measures the error structure instead of assuming one.

  \item \textbf{Never evaluate a staggered rollout with a single post dummy --- and get the reason
        right.} With growing effects, two-way fixed effects lands at \eur{\nbEightTwfeNaive} against
        a true \eur{\nbEightTwfeTrue}, below every true effect in the panel. The
        \citeauthor{goodmanbacon2021} weights on the two $2 \times 2$s are both \emph{positive}
        (\nbEightBaconWClean{} and \nbEightBaconWForbidden, \cref{tab:did:bacon}); the damage rides
        in through the forbidden comparison's \emph{estimate} ---
        \eur{\nbEightBaconDidForbidden} against \eur{\nbEightBaconDidClean} --- which measures a late
        cohort against an already-treated one. The negative weights are \citeauthor{dechaisemartin2020}'s,
        and they sit on the \emph{cells}: \nbEightBaconNegCells{} of \nbEightBaconTreatedCells{}
        cohort-weeks, the largest effects in the panel \citep{goodmanbacon2021, dechaisemartin2020}.
        \Citeauthor{callaway2021}'s discipline (\cref{eq:did:cs}: only not-yet-treated controls)
        recovers \eur{\nbEightCsAtt} $\pm$ \eur{\nbEightCsSe} (\cref{tab:did:staggered}). And
        \textbf{keep a never-treated holdout to the very end}: without one, $\mathcal C(t)$ empties
        and the later waves become unmeasurable.

  \item \textbf{Significant is not the same as scalable.} The effect is real
        ($p \le \nbEightPermP$) and worth \eur{\nbEightAnnualMeanM}\,m a year \emph{if the chain
        responds like the pilot}. The estimand is the ATT --- the effect \emph{on the pilot stores}
        --- and \cref{tab:did:stress} shows the two axes interacting: at the vendor's
        \eur{\nbEightAppCost} the call survives even half transfer ($\nbEightPHalfAppCost$), but the
        transfer share sets how much the chain can afford to pay. The deliverable is the break-even
        band from the \emph{calibrated} posterior
        (\eur{\nbEightCapHalfFe}--\eur{\nbEightCapFullFe} per store-week; the condemned pooled model
        would have quoted \eur{\nbEightCapHalf}--\eur{\nbEightCapFull} and rejected quotes it should
        accept), not the base-case annual figure --- and the cheap next move is a second wave in
        deliberately non-pilot-like stores.
\end{enumerate}

\notebooknote{%
  This chapter is \code{notebooks/08\_rollout\_did.ipynb}, which runs in a fast teaching mode
  (\code{CMP\_FAST=1}, short chains) and a full mode (\code{CMP\_FAST=0}); every number here comes
  from the full run, emitted by \code{cmp.report} and injected as a macro. The simulator is
  \code{cmp.dgp.did\_rollout}, the classical estimators are \code{cmp.classical} (\code{did\_2x2},
  \code{event\_study}, \code{ols} with \code{cov="cluster"}), the Bayesian $2 \times 2$ is CausalPy
  \citep{causalpy}, and the seeds are fixed in the notebook. \Citeauthor{card1994}'s minimum-wage
  study of fast-food employment in New Jersey and Pennsylvania \citep{card1994} remains the
  canonical public panel on which to practise the method, and \citet{roth2023} is the current map of
  the staggered-adoption literature this chapter's \cref{sec:did:wrong} only samples.}
